\hspace{3mm}

\centerline{\bf \Huge Оценка+пример}

\centerline{\bf \large Базовый уровень}

\begin{flushright}
        \scriptsize{}
\end{flushright}

\textbf{Оценка плюс пример} --- это специальное математическое рассуждение, которое применяется в задачах, где требуется найти наибольше или наименьшее значение чего либо. Сначала доказывается, что требуемая величина как-то ограничена сверху или снизу (так называемая \textit{оценка}) а затем приводят пример, когда эта самая оценка достигается.  

\problem{0.1} 
Какое наибольшее число трёхклеточных уголков можно вырезать из клетчатого квадрата $8 \times 8$?

\problem{0.2} 
Каким наименьшим числом монет в $3$ и $5$ копеек можно набрать сумму $37$ копеек?

\problem{1}
Какое наибольшее число трёхклеточных уголков можно вырезать из клетчатого прямоугольника $5 \times 7$?

\problem{2}
Какое наименьшее число выстрелов в игре "Морской бой" на доске 7*7 нужно сделать, чтобы наверняка ранить четырехпалубный корабль (четырехпалубный корабль состоит из четырех клеток, расположенных в один ряд)?

\problem{3}
На шахматной доске стоят 10 шахматных фигур --- слоны и ладьи --- не бьющие друг друга. Какое наименьшее количество слонов может быть среди них? 

\begin{wrapfigure}{r}{0.3\linewidth}
    \vspace{-7mm}
    \includegraphics[width=\linewidth]{Летние Сборы/Сборы Элиста 2024/АА/Оценка+пример.png}
\end{wrapfigure}


\problem{4}
План дома Германа имеет форму равностороннего треугольника со стороной 100 м. Он разделён на 100 треугольных залов (рис.). Все стены залов имеют одинаковую длину — 10 метров. В середине каждой стены между залами сделаны двери. В каждом зале находится по одному ШкибидиКоину. Какое наибольшее количество ШкбидиКоинов сможет собрать Вова, если в каждой комнате можно побывать ровно один раз?

\problem{5}
В турнире участвуют 100 борцов, все разной силы. Более сильный всегда побеждает более слабого. Борцы разбились на пары и провели поединки. Затем разбились на пары по-другому и снова провели поединки. Призы получили те, кто выиграл оба поединка. Каково наименьшее возможное количество призѐров?


\newpage

\hspace{3mm}

\centerline{\bf \Huge Оценка+пример}

\centerline{\bf \large Продвинутый уровень}

\begin{flushright}
        \scriptsize{}
\end{flushright}

\problem{1} 
В пруд пустили 30 лососей, которые постепенно поедают друг друга. Лосось считается сытым, если он съел не менее трёх лососей (сытых или голодных). Какое наибольшее число лососей может насытиться?

\problem{2} 
Какое наименьшее количество множителей требуется вычеркнуть из числа 99! так, чтобы произведение оставшихся множителей оканчивалось на 2?

\problem{3} 
 Каково наибольшее количество последовательных натуральных чисел, у каждого из которых ровно четыре натуральных делителя (включая 1 и само число)?

\problem{4}
Найдите наибольшее натуральное число, которое невозможно представить в виде суммы двух составных чисел.

\problem{5} 
Изначально на экране компьютера Дамира --- какое-то простое число. У Дамира есть кнопка, которая заменяет число на число, полученное из предыдущего прибавлением его последней цифры, увеличенной на 1. Каждый раз, когда у Дамира появляется простое число он получает 1 ЭлМШкоин, а как только появляется составное число у Дамира сгорает видоекарта и игра заканчивается. Какое наибольшее количество ЭлМШкоинов может заработать лучший ученик Калмыкии?